\documentclass[11pt,a4paper]{article}

\usepackage[utf8]{inputenc}
\usepackage[T1]{fontenc}
\usepackage[margin=2.5cm]{geometry}
\usepackage{amssymb}
\usepackage{parskip}
\usepackage{microtype}
\usepackage{hyperref}
\usepackage{xurl}
\usepackage{csquotes}
\usepackage{enumitem}
\usepackage{fancyvrb}

\hypersetup{
  colorlinks=true,
  linkcolor=black,
  urlcolor=blue,
  citecolor=blue,
  breaklinks=true
}

\emergencystretch=3em
\sloppy

\setlist[itemize]{label=$\blacktriangleright$}

\newcommand{\code}[1]{\texttt{#1}}
\newcommand{\pkg}[1]{\textsf{#1}}

\title{Response to the Second-Round Review\\
\large \textit{BayesianDEB}: A Bayesian Framework for Dynamic Energy
Budget Modelling in R}
\author{Branimir K.\ Hackenberger\\ \small on behalf of the authors}
\date{2026-06-29}

\begin{document}
\maketitle

We are grateful to the editor and reviewer for reading the manuscript so
closely a second time. Each comment from the review of \textbf{22 May
2026} is quoted below in italics, with our reply underneath and the files
we touched to address it. Everything here refers to version
\textbf{0.2.1} of the package; the replication archive is
\code{BayesianDEB\_replication.zip} and the revised manuscript is
\code{bayesiandeb.tex}.

In short:
\begin{itemize}[nosep]
\item a single master replication script (\code{replicate\_all.R})
  rebuilds every manuscript number in about a minute from cached draws;
\item every code listing in the manuscript now runs, and prints its
  output, in the replication scripts;
\item all manuscript outputs come from the shipped cache under a fixed
  seed, so they reproduce bit for bit;
\item the diagnostics \code{print} and \code{plot} methods are now
  compact;
\item the \code{bdeb\_diagnose} examples have moved from
  \code{\textbackslash dontrun\{\}} to a \pkg{cmdstanr}-gated
  \code{\textbackslash donttest\{\}} block;
\item the dose-response figure now draws the reference lines its caption
  describes;
\item the README installation instructions no longer pull the stale CRAN
  version.
\end{itemize}

\section{Package version must match the CRAN version}

\begin{displayquote}
``The submitted package version is not in line with the package
published on CRAN. We require that both versions match.''
\end{displayquote}

\textbf{Response.} When the reviewer wrote, CRAN still held 0.1.4 while
the manuscript and this submission describe 0.2.1. That gap is now
closed: 0.2.1 went up on CRAN on 2026-06-17, so the published package and
the manuscript describe the same release. The version bump is documented
in \code{cran-comments.md}, and \code{DESCRIPTION} declares
\code{Version: 0.2.1}. \verb|R CMD check --as-cran| reports 0 errors,
0 warnings and a single note (the informational
\pkg{cmdstanr}-in-\code{Additional\_repositories} line, discussed below).

\section{README should include cmdstanr installation instructions}

\begin{displayquote}
``The README file should also include instructions for the installation
of cmdstanr.''
\end{displayquote}

\textbf{Response.} Added. The Installation section of \code{README.md}
now spells out the two steps: first install \pkg{cmdstanr} from the Stan
r-universe, then build the CmdStan toolchain once,
\begin{Verbatim}[fontsize=\small]
install.packages("cmdstanr",
  repos = c("https://stan-dev.r-universe.dev", getOption("repos")))
cmdstanr::install_cmdstan()
\end{Verbatim}
\noindent and the text notes that \code{bdeb\_fit()} checks for the
toolchain at run time.

\section{A single replication script that runs within one hour}

\begin{displayquote}
``Journal of Statistical Software requires that a single replication
script that can be run within one hour is provided. \ldots\ parts of the
results require that the scripts in \code{sbc} are run and these scripts
are supposed to take more than 20 hours to finish. We would thus ask
that the authors arrange their replication scripts into a smaller number
of R files and that they also work on reducing the computational time
needed to replicate results.''
\end{displayquote}

\textbf{Response.} There is now a single master script,
\code{replication/replicate\_all.R}, that rebuilds every figure, table
and printed number in the manuscript from one command
(\verb|Rscript replicate_all.R|). Out of the box it runs no MCMC at all:
it reads the archived posterior draws in \code{outputs/} and \code{sbc/}
and regenerates everything in roughly a minute on a laptop, comfortably
under the one-hour limit. The lengthy simulation-based-calibration runs
stay out of this path. Their rank matrices are shipped
(\code{sbc/*.rds}) and reproduce Figures~3--4 and Tables~4--5 on their
own; recomputing them from scratch takes anywhere from 45~min to 12~h per
model and is described separately in \code{sbc/README.md}, outside the
replication budget. Anyone who does want to refit can set
\code{BDEB\_RECOMPUTE=true BDEB\_MODE=full} for a complete run
($\sim$3~h) or \code{BDEB\_MODE=lite} for a quick one ($\sim$30~min).

\textbf{Reference.} \code{replication/replicate\_all.R},
\code{replication/README.md}, \code{replication/sbc/README.md}.

\section{\code{example("bdeb\_diagnose")} should be executable
(\code{\textbackslash donttest})}

\begin{displayquote}
``\verb|example("bdeb_diagnose")| shows that this example is in
\code{\textbackslash dontrun}. This should, in principle, be executable
code and we expect it to be in \code{donttest} instead so as to allow
its execution using \code{example}, similarly to what has been
implemented for \code{bdeb\_fit}.''
\end{displayquote}

\textbf{Response.} Done. The examples for \code{bdeb\_diagnose()}, the
\code{bdeb\_diagnostics} \code{print}/\code{summary}/\code{plot} methods,
\code{bdeb\_loo()}, \code{bdeb\_derived()} and the deprecated
\code{bdeb\_summary()} are now wrapped in the same \pkg{cmdstanr}-gated
\code{\textbackslash donttest\{\}} block we already use for
\code{bdeb\_fit()}:
\begin{Verbatim}[fontsize=\small]
\donttest{
if (requireNamespace("cmdstanr", quietly = TRUE) &&
    nzchar(tryCatch(cmdstanr::cmdstan_path(), error = function(e) ""))) {
  data(eisenia_growth)
  dat <- bdeb_data(growth = eisenia_growth[eisenia_growth$id == 1, ])
  fit <- bdeb_fit(bdeb_model(dat, type = "individual"),
                  chains = 2, iter_warmup = 200, iter_sampling = 200,
                  refresh = 0)
  print(bdeb_diagnose(fit))
}
}
\end{Verbatim}
\noindent so \verb|example("bdeb_diagnose")| runs whenever a CmdStan
toolchain is available. None of the exported functions rely on
\code{\textbackslash dontrun\{\}} any more.

\textbf{Reference.} \code{R/diagnostics.R}, \code{R/fit.R},
\code{R/ppc.R}, \code{R/debtox.R}.

\section{The \code{plot} method for \code{bdeb\_diagnostics} could be
more readable}

\begin{displayquote}
``We believe that the plot method for \code{bdeb\_diagnostics} could be
improved to allow a shorter and more readable display on screen.''
\end{displayquote}

\textbf{Response.} \code{plot.bdeb\_diagnostics()} no longer crowds the
display with per-time-point latent states (\code{x\_sol[i,j]},
\code{L\_hat[i]}, and so on). By default it shows only the scalar model
parameters, so the R-hat and ESS panels stay short and easy to read
rather than running to dozens of latent rows. Passing \code{full = TRUE}
brings back the complete plot, and the subtitle notes how many latent
rows were left out. This follows what \code{print.bdeb\_diagnostics()}
already does.

\textbf{Reference.} \code{R/diagnostics.R}
(\code{plot.bdeb\_diagnostics}); \code{NEWS.md}.

\section{Replication output differs from the article (LOO comparison)}

\begin{displayquote}
``Results obtained with the replication material (\code{full} version)
are still different from the results shown in the article. For instance,
on page~13, the \code{loo\_compare} for the lognormal model differs.''
\end{displayquote}

\textbf{Response.} This comparison really is sensitive to Monte-Carlo
noise on this individual. With only $n = 13$ observations, a handful of
Pareto-$\hat{k}$ values climb above $0.7$, so the importance-sampling
ELPD difference moves by a fraction of its standard error from run to
run; across our reruns it landed around $-1$ every time, always with SE
$\approx 1.3$ and always within one standard error of zero. We made
three changes:
\begin{itemize}[nosep]
\item fixed a global seed (\code{seed = 42}) in every fit, so the shipped
  result is deterministic;
\item had \code{replicate\_all.R} reproduce the manuscript numbers from
  the cached draws, so the printed $\Delta$ELPD matches the article
  exactly;
\item rewrote the conclusion as a qualitative one: since
  $|\Delta\mathrm{ELPD}|$ is smaller than its standard error, neither
  observation model wins clearly, and the precise difference is best read
  as indicative.
\end{itemize}
The manuscript prints \code{lognormal -1.0} (SE~1.3) and says outright
that the difference sits within its standard error.

\textbf{Reference.} \code{replication/00\_setup.R} (seed),
\code{replication/01\_illustrations.R}, manuscript Section~5.1.

\section{README installation pulls the old CRAN version (0.1.4)}

\begin{displayquote}
``The instructions in the \code{README.md} suggest installing with a
\code{repos} option that includes CRAN, which installs version 0.1.4 and
not the submitted package version 0.2.0.''
\end{displayquote}

\textbf{Response.} We have split the installation instructions.
\pkg{BayesianDEB} is now installed with a plain
\verb|install.packages("BayesianDEB")|, which now resolves to 0.2.1 from
CRAN, while \pkg{cmdstanr} is installed separately from the Stan
r-universe (see Section~2). The earlier combined
\code{repos = c("https://stan-dev.r-universe.dev", getOption("repos"))}
call, which could have masked the CRAN copy of \pkg{BayesianDEB}, is gone
from the README.

\textbf{Reference.} \code{README.md} (Installation section).

\section{Replication material is incomplete (missing manuscript code)}

\begin{displayquote}
``Some code included in the manuscript which should be executable code is
not contained \ldots\ E.g., on page~7,
\verb|prior_species("Daphnia_magna", type = "debtox")|.''
\end{displayquote}

\textbf{Response.} Every code listing in the manuscript is now run in the
replication scripts. \verb|prior_species("Daphnia_magna", type = "debtox")|,
for example, is called and printed in
\code{replication/01\_illustrations.R} next to the
\verb|prior_species("Eisenia_fetida")| listing. We went back through the
manuscript one listing at a time and checked that each appears in a
replication script.

\textbf{Reference.} \code{replication/01\_illustrations.R}.

\section{\code{bdeb\_diagnose(fit)} output differs and is too long}

\begin{displayquote}
``Some results obtained differ. E.g., compared to page~12,
\code{bdeb\_diagnose(fit)} reports a different number of divergent
transitions and prints a very long table including \code{x\_sol[i,j]}
and \code{L\_hat[i]} rows.''
\end{displayquote}

\textbf{Response.} Two things address this:
\begin{itemize}[nosep]
\item \textbf{Reproducibility.} With the fixed seed and cached draws
  (Section~6), \code{replicate\_all.R} reproduces the page-12 output
  exactly; the manuscript now reports a single divergent transition
  (1 of 4000 draws, $<0.1\%$) and comments on it.
\item \textbf{Length.} \code{print.bdeb\_diagnostics()} now defaults to
  the scalar model parameters and hides the per-time-point latent
  states, so the p.~12 output is compact (eight parameter rows, no
  \code{x\_sol}/\code{L\_hat}). The full table remains available through
  \code{print(x, full = TRUE)} or \code{summary(x)\$table}.
\end{itemize}

\textbf{Reference.} \code{R/diagnostics.R},
\code{replication/01\_illustrations.R}, manuscript Section~5.1.

\section{\code{bdeb\_derived()} result is created but never printed}

\begin{displayquote}
``The replication material creates \code{der <- bdeb\_derived(fit, ...)}
but does not print it, so it is not easy to compare with the
manuscript.''
\end{displayquote}

\textbf{Response.} Fixed. \code{replication/01\_illustrations.R} now
prints the derived-quantity table right after it is built, laid out as in
the manuscript:
\begin{Verbatim}[fontsize=\small]
der <- bdeb_derived(fit,
  quantities = c("L_m", "L_inf", "k_M", "g", "growth_rate"), f = 1.0)
print(as.data.frame(summarise_draws(der, "mean", "sd",
  "q5" = ~quantile(.x, 0.05), "q95" = ~quantile(.x, 0.95))), digits = 3)
\end{Verbatim}

\textbf{Reference.} \code{replication/01\_illustrations.R}.

\section{LOO triggers Pareto-$\hat{k}$ warnings that are not commented
on}

\begin{displayquote}
``Warnings are triggered (\code{Some Pareto k diagnostic values are too
high}) where one might also point them out and comment on them.''
\end{displayquote}

\textbf{Response.} These warnings are now flagged and explained rather
than left to surprise the reader. In
\code{replication/01\_illustrations.R} a comment ahead of the
\code{loo\_compare()} call points out that, at $n = 13$, a few
$\hat{k} > 0.7$ are to be expected and the importance-sampling ELPD is
therefore only indicative. The manuscript makes the same point in the
text of Section~5.1: it reports the Pareto-$\hat{k}$ warning, explains
what it means (the approximation is unreliable at a few points), and uses
it to justify reading the comparison qualitatively.

\textbf{Reference.} \code{replication/01\_illustrations.R}, manuscript
Section~5.1.

\section{Figure caption refers to lines that are not in the figure}

\begin{displayquote}
``The caption of Figure~2 discusses something dashed red and dotted
green. However, it is unclear to what in the figure one refers.''
\end{displayquote}

\textbf{Response.} \code{plot\_dose\_response()} now draws the lines the
caption promises: a dashed red vertical line at the posterior-median
EC$_{50}$ and a dotted green one at the posterior-median NEC, with a
matching legend subtitle. Non-finite draws are discarded and the view is
clipped with \code{coord\_cartesian()}, so degenerate draws no longer
leave vertical artefacts. The Figure~2 caption and the rendered figure
now agree.

\textbf{Reference.} \code{R/debtox.R} (\code{plot\_dose\_response}),
manuscript Figure~2 caption.

\section{Unclear which part of the replication reproduces Figures 3 and
4}

\begin{displayquote}
``It is unclear what part of the replication material reproduces
Figures~3 and~4.''
\end{displayquote}

\textbf{Response.} \code{replication/README.md} now carries a
\emph{Figure $\rightarrow$ script} table that ties each manuscript figure
to the script and cached object behind it. For the two in question:
Figure~3 (\code{fig\_sbc\_ranks.pdf}) comes from \code{02\_validation.R}
via \code{sbc/sbc\_results.rds}, and Figure~4
(\code{fig\_sbc\_hierarchical.pdf}) comes from \code{02\_validation.R}
via \code{sbc/sbc\_hierarchical\_results.rds}. The same table covers the
remaining figures and the matching manuscript sections.

\textbf{Reference.} \code{replication/README.md} (Figure $\rightarrow$
script and Section $\rightarrow$ script tables).

\vspace{1em}
\noindent We thank the editor and reviewer once more for feedback that
has made the package and its replication material better.

\vspace{1em}
\noindent Branimir K.\ Hackenberger\\
\noindent on behalf of the authors\\
\noindent 2026-06-29

\end{document}
